% TOP_PROBLEM: {"id": 20000908, "title": "Higher additive energy and quantitative stability of the Holder defect", "category_id": 2, "category": {"id": 2, "name": "combinatorics", "display_name": "Combinatorics", "description": "Counting problems, graph theory, discrete structures.", "slug": "combinatorics", "order_index": 2, "created_at": "2026-07-31T15:26:25.671Z"}, "status": "partially_solved", "problem_number": "AIM-COMBINATORICS-0033", "set_id": 14, "statusQualification": "Includes the source discussion correction |A|^epsilon and permits constants to depend on n and epsilon; repairs the tuple-index transcription."}
% FIRST_POSTED: 2026-10-01
% ENTRY_KIND: statement_only
% UnsolvedMath ID 20000908; source catalogue code AIM-COMBINATORICS-0033
% !TeX program = lualatex
% Definition and sources only; this file is not an LLM solution attempt.
\documentclass[11pt,a4paper]{article}
\usepackage[margin=25mm]{geometry}
\usepackage{fontspec}
\setmainfont{lmroman10-regular.otf}[BoldFont=lmroman10-bold.otf,
 ItalicFont=lmroman10-italic.otf,BoldItalicFont=lmroman10-bolditalic.otf]
\usepackage{amsmath,amssymb,mathtools}
\usepackage{unicode-math}
\setmathfont{latinmodern-math.otf}
\AtBeginDocument{\let\setminus\smallsetminus}
\usepackage{xurl,hyperref}
\hypersetup{colorlinks=true,urlcolor=blue}
\setcounter{secnumdepth}{0}
\setlength{\parindent}{0pt}
\setlength{\parskip}{5pt}
\setlength{\emergencystretch}{3em}
\begin{document}
\section*{Higher additive energy and quantitative stability of the Holder defect}\label{20000908}
{\small UnsolvedMath ID 20000908 · AIM-COMBINATORICS-0033 · Combinatorics · AIM Workshop Problem Lists}
\subsection{Definitions and mathematical statement}
Let $A$ be a finite nonempty subset of an abelian group $G$.
For $j\geq2$, define its higher additive energies and their normalizations by
\[
 E_{2j}(A)=|\{(a_1,\ldots,a_{2j})\in A^{2j}:
       a_1+\cdots+a_j=a_{j+1}+\cdots+a_{2j}\}|,
 \qquad \rho_{2j}(A)=E_{2j}(A)/|A|^j.
\]
These are counts of ordered tuples, allowing repetitions.
For an integer $m\geq1$, $mA$ denotes the set of sums of $m$ elements of $A$, again allowing repetitions.

Fix an integer $n\geq2$ and a real $\varepsilon\in(0,1)$.
Does there exist a pair of positive constants $c,c'$, depending only on $n,\varepsilon$, such that the following assertion holds for every $G,A$ and every $K>1$?
If
\[
 \rho_{2n+2}(A)^{(n-1)/n}<K\rho_{2n}(A),
\]
then some nonempty $H\subseteq(n-1)A$ satisfies
\[
 |H|\geq\frac{\rho_{2n}(A)^{1/(n-1)}}{K^c},
 \qquad E_4(H)\geq\frac{|H|^3}{|A|^{\varepsilon}K^{c'}}.
\]
Thus near equality in the comparison between successive normalized energies should produce a relatively large set with substantial ordinary additive energy.
The constants must not depend on $A$, its ambient group, or $K$.
The factor $|A|^{\varepsilon}$ in the energy denominator is retained: the source discussion explicitly corrects a formulation that had omitted this loss.
No claim that the extracted set is a subgroup is part of the question.
\subsection{Short English statement}
Does slow growth between successive higher additive energies force a large, energetically structured subset of an iterated sumset?
\subsection{Sources}
\begin{itemize}
\item[S1] ULAM.AI and the UnsolvedMath Contributors, supplied problems.json, record 20000908 (AIM-COMBINATORICS-0033), AIM Workshop Problem Lists. Catalogue identity and source transcription; CC BY 4.0.
\url{https://huggingface.co/datasets/ulamai/UnsolvedMath}.
\item[S2] American Institute of Mathematics, Additive combinatorics and its applications, item 2.6. Original workshop question underlying AIM-COMBINATORICS-0033.
\url{http://aimpl.org/addcombapp/2/}.
\end{itemize}
\end{document}
